%=====================================================================
\section{M15: Profiling und Parallelisierung}\label{sec:parallel}
%=====================================================================

\subsection{Ziel}
Eine \emph{einzelne} Rechnung (nicht nur ein Scan) nutzt alle Kerne bzw.\ die GPU gut aus. Ergebnisse sind
deterministisch und unabhängig von der Threadzahl.

\subsection{Schritt 1: Messen}
\begin{itemize}[nosep]
  \item \verb|fg::prof|: Scope-Timer (\verb|FG_PROF_SCOPE("rcwa.eig")|), threadsichere Aggregation (Summe, Anzahl, Max
        je Name), Ausgabe als JSON (\verb|fgsolve --profile|, Python \verb|res.diagnostics["timings"]|). Kosten
        abgeschaltet $<1\,\%$.
  \item Profile für V3, V9, V10 ($N=10$), V13 und das CAD-Beispiel auf dem Zielsystem (CPU und GPU) in
        \verb|docs/benchmarks.md| (\enquote{Profil vor M15}). Die folgenden Schritte werden nach diesem Profil priorisiert;
        was nicht messbar hilft, entfällt.
\end{itemize}

\subsection{Schritt 2: Threadmodell}
Drei Ebenen: Scanpunkte (bestehend, \verb|solveScan|), Schichten/Scheiben einer Rechnung (neu), BLAS/Eigen
innerhalb einer Matrixoperation. Konfiguration \verb|"threads"| mit den Schlüsseln \verb|scan|, \verb|slices| (je Zahl oder \verb|auto|) und \verb|blas| (Standard 1),
Umgebungsvariable \verb|FG_NUM_THREADS|. Regel gegen Überbelegung: Produkt der Ebenen $\le$ Kernzahl; \verb|auto|
verteilt zuerst auf die Scanebene (wenn Scanpunkte $\ge$ Kerne), sonst auf die Scheibenebene. OpenMP-Tasks;
verschachtelte Parallelität explizit (\verb|omp_set_max_active_levels(2)|), Eigen über \verb|Eigen::setNbThreads|.

\subsection{Schritt 3: Parallele Scheiben-S-Matrizen und Baumreduktion}
\begin{itemize}[nosep]
  \item Die S-Matrizen aller Schichten und Scheiben (RCWA-Scheiben, RK-Stabilisierungsscheiben) sind unabhängig und werden
        parallel berechnet; Ablage nach Index.
  \item Das Sternprodukt ist assoziativ. Verkettung als balancierter Binärbaum mit \emph{fester} Form (Paare $(0,1),(2,3),\dots$,
        dann rekursiv), unabhängig von der Threadzahl $\Rightarrow$ bitgleiche Ergebnisse für jede Threadzahl.
        Abweichung zur sequentiellen Verkettung $\le10^{-12}$ relativ (Test). Option \verb!"reduction": "tree" | "sequential"!.
  \item Nahfeld (\verb|storeInternal|) braucht $S_L$ und $S_U$ an jeder inneren Grenzfläche: zunächst sequentiell
        (Präfix/Suffix); parallele Präfixsumme nur, wenn das Profil es verlangt.
  \item RK: Stufenauswertungen $Q(y+c_ih)$ einer Scheibe parallel; Scheiben untereinander parallel.
\end{itemize}

\subsection{Schritt 4: Geometrie und Koeffizienten}
Raster von Algorithmus~NF, Konturquadratur (je Thread eigenes Teilgitter, Summation in fester Reihenfolge),
Profil-Schnittpunkte und CAD-Schnitte je Höhe parallel. Für Nurbforge nur die als threadsicher dokumentierten
\verb|const|-Abfragen verwenden (M14).

\subsection{Schritt 5: GPU-residente Ketten}
Das CUDA-Backend arbeitet bisher über die Host-API (jede Operation kopiert Ein- und Ausgaben). Bei FP32 (M16) wird der
PCIe-Transfer sonst zum Engpass (Beispiel $n=3362$: eine Matrix 90--180\,MB, CGEMM auf der RTX~3090 nur wenige ms).
Daher:
\begin{itemize}[nosep]
  \item \verb|linalg::DeviceMatrix| (RAII-Handle) und Operationen \verb|gemm|, \verb|solve|, \verb|eig|, \verb|add|,
        \verb|scaleRows| auf Handles; \verb|upload|/\verb|download| explizit.
  \item RCWA-Scheiben-S und Sternproduktkette laufen vollständig auf dem Gerät; nur das Gesamt-$S$ (bzw.\ die für das
        Nahfeld nötigen Matrizen) wird zurückkopiert. Mehrere Streams für unabhängige Scheiben.
  \item Kontrolle mit Nsight Systems: Anteil der Transferzeit $<10\,\%$ bei 2D, $N\ge10$.
\end{itemize}

\paragraph{Umsetzung (v2.7).} Die Baumform wird ohne vollständige Ablage aller Scheiben erreicht: ausgerichtete
Blöcke von $2^m\ge$ Threads Elementen werden parallel paarweise reduziert, vollständige Teilbäume wie ein
Binärzähler zusammengefasst und der Rest von rechts gefaltet; das ergibt dieselbe Form wie die stufenweise
Paarbildung für jede Blockgröße (Test bitgleich), Speicher $O(\text{Threads}+\log K)$. Standard bleibt
\verb!"reduction": "sequential"! (bitgleich zu v0.10); die parallele Berechnung der Scheiben selbst ändert keine Bits und
ist immer aktiv. \verb|auto| gibt der Scanebene $\min(\text{Scanpunkte},\text{Kerne})$ Threads und der Scheibenebene den
Rest. Innere Schleifen (Raster von Algorithmus~NF zeilenweise, Gauß-Glättung, Konturstücke je Segment,
Tschebyschow-Knoten) laufen nur parallel, wenn der Aufrufer selbst sequentiell ist; Summationsreihenfolgen bleiben
erhalten (bitgleich). Ohne OpenMP verteilt \verb|parallelFor| auf \verb|std::thread|; der TSAN-Job baut so, weil
libgomp nicht instrumentiert ist. MSVC: OpenMP~2.0 genügt (\verb|omp_set_nested| statt \verb|omp_set_max_active_levels|).
GPU: \verb|DeviceMatrix| mit Speicherpool je Backend (kein \verb|cudaFree| im Lauf); ein Stream je Thread, eine Matrix
wird erst nach \verb|synchronize()| des erzeugenden Threads weitergegeben, beim Freigeben durch einen anderen Thread
wartet dieser auf seinen Stream. Die residenten Operationen haben eine Host-Referenz (Eigen), gegen die der
Gerätepfad auch ohne GPU geprüft wird. Höchstens \verb|FG_GPU_STREAMS| (Standard 4) Scheiben gleichzeitig auf dem Gerät.

\subsection{Tests und Abnahme (M15)}
\begin{itemize}[nosep]
  \item Determinismus: Threadzahl 1, 2, 4 liefern bitgleiche Ergebnisse (alle Validierungsfälle, kleine $N$).
  \item Baum = sequentiell auf $10^{-12}$; GPU-residente Kette = Host-Kette auf $10^{-10}$ (GPU-Rechner).
  \item CI-Job mit ThreadSanitizer (Unit-Tests, kleine Größen) ohne Befund.
  \item Leistungsbericht in \verb|docs/benchmarks.md| (Zielsystem und RTX~3090): Laufzeit einer 2D-Rechnung mit
        $\ge32$ Scheiben vor/nach, parallele Effizienz (Zielwert $\ge60\,\%$ auf 16 Kernen), Transferanteil GPU.
\end{itemize}

%=====================================================================
\section{M16: Wählbare Genauigkeit (FP64 / gemischt / FP32)}\label{sec:precision}
%=====================================================================

\subsection{Motivation und Grundsatz}
Auf der RTX~3090 (GA102, \verb|sm_86|) und der RTX 6000 Ada (\verb|sm_89|) beträgt die FP64-Rate etwa 1/64 der
FP32-Rate. Die $O(n^3)$-Operationen (GEMM, LU, Eigenwerte) dominieren die Laufzeit. Die Genauigkeit wird deshalb
\textbf{am Backend} gewählt: Matrizen bleiben im Physikcode \verb|MatC| (FP64); das Backend wandelt an seiner Grenze
($O(n^2)$) um und rechnet die $O(n^3)$-Operation in der gewählten Genauigkeit. Der Physikcode wird nicht templatisiert.

\subsection{Modi}
\begin{center}
\small
\begin{tabular}{p{4.6cm}p{2.6cm}p{3.6cm}p{3.0cm}}
\toprule
Operation & \verb|fp64| (Standard) & \verb|mixed| & \verb|fp32| \\
\midrule
Toeplitz, $Q$, $M$, Koeffizienten, Normalenfeld & FP64 & FP64 & FP64 \\
GEMM & Z & C & C \\
LU / Lösen & Z & C-Faktorisierung + FP64-Nachiteration (IRS) & C \\
Eigenwerte (RCWA) & Z & C & C \\
Exponentiale, stabile $S$-Formeln, Effizienzen & FP64 & FP64 & FP64 \\
\bottomrule
\end{tabular}
\end{center}
GPU: \verb|cublasCgemm| (Math-Mode ausdrücklich \verb|CUBLAS_DEFAULT_MATH|, \textbf{kein TF32}; zusätzlich
\verb|NVIDIA_TF32_OVERRIDE=0| in der Doku), \verb|cusolverDnXgetrf|/\verb|Xgetrs| und \verb|cusolverDnXgeev| mit
\verb|CUDA_C_32F|, gemischtes Lösen mit der IRS-API (\verb|cusolverDnZCgesv| bzw.\ \verb|cusolverDnIRSXgesv|).
CPU: Eigen \verb|MatrixXcf| (bzw.\ LAPACK \verb|cgetrf|/\verb|cgeev|), Nachiteration selbst implementiert:
$x\leftarrow x+A^{-1}_{32}(b-Ax)$ mit Residuum in FP64, höchstens 10 Schritte, Abbruch bei
$\|b-Ax\|\le\epsilon_{64}\|A\|\|x\|\cdot10$, sonst Rückfall auf FP64-LU (Diagnose).

\subsection{Angepasste Parameter}
\begin{itemize}[nosep]
  \item Scheibenkriterium \eqref{eq:slicecrit}: Wachstumsgrenze $G_{\max}=\epsilon^{-1/2}$ der Arbeitsgenauigkeit der
        Propagation. FP64: $6.7\cdot10^{7}\approx10^8$ (wie bisher), FP32/mixed: $\approx4\cdot10^{3}$, also
        $\ln G_{\max}\approx8.3$ statt 18.4 $\Rightarrow$ mehr, dünnere Scheiben. \verb|growthLimit| bekommt den Standard
        \verb|auto| (abhängig vom Modus); ein expliziter Wert gilt weiter.
  \item RK: \verb|rk.tol| wird in FP32 auf $\ge10^{-5}$ angehoben (Warnung).
  \item Eigenwertklassifikation: Toleranzen relativ zu $\max|\lambda|$ mit $\epsilon_{32}$ skalieren.
\end{itemize}

\subsection{Schutzmechanismen}
Diagnosen je Rechnung: Energiebilanz (verlustfrei), Konditionszahlen ($T_{11}$, $G$), Zahl der Nachiterationen.
Option \verb|"precision_fallback": true| (Standard in \verb|fp32|): Rechnung automatisch in FP64 wiederholen, wenn bei
verlustfreien Strukturen $|\sum e-1|>10^{-3}$ oder $\operatorname{cond}(T_{11})>10^{5}$; Warnung im Ergebnis.

\subsection{Standardwahl (Entscheidung des Projektinhabers)}
\verb|"precision": "auto"| ist der Standard und bedeutet: CPU-Pfad immer \verb|fp64|; auf der GPU \verb|fp32| mit
\verb|precision_fallback| auf einer GeForce RTX~3090 (Anwenderinstallation Windows~11), sonst \verb|fp64|
(Server, RTX 6000 Ada). Erkennung über \verb|cudaGetDeviceProperties| (Name und Compute Capability 8.6); die Regel steht
in einer Tabelle \verb|src/cuda/device_policy.cpp|, damit weitere Karten ergänzt werden können. Überschreibbar mit
\verb|FG_PRECISION| und in der Konfiguration. Das Ergebnis meldet die tatsächlich verwendete Genauigkeit
(\verb|diagnostics.precision|) und ob ein Rückfall auf FP64 stattfand.

\subsection{Build}
\verb|CMAKE_CUDA_ARCHITECTURES| \verb|"86;89"| in den CUDA-Presets (RTX 3090 und RTX 6000 Ada); Preset
\verb|win-cuda| ebenfalls \verb|86|. Speichertabelle (Abschnitt Architektur) um FP32 und 24\,GB (3090) ergänzen.

\paragraph{Umsetzung (v2.8).} \verb|linalg::Precision| wird je Thread gesetzt; in FP32/Mixed läuft jede
$O(n^3)$-Operation über ein Backend (GPU über der Schwelle, sonst ein CPU-Backend in einfacher Genauigkeit: Eigen
\verb|MatrixXcf|, eigene Nachiteration). Damit ist die Studie auch ohne GPU möglich. Auf der GPU tragen die
Gerätematrizen den Elementtyp (Z oder C); die Host-API (\verb|gemm|/\verb|solve|/\verb|eig| auf \verb|MatC|) ist auf
den residenten Operationen aufgebaut. FP32 läuft resident (Umwandlung beim Hoch-/Herunterladen); Mixed nutzt die
Host-API, weil die Typumwandlung auf dem Gerät Kernel bräuchte, und löst mit \verb|cusolverDnZCgesv|. Die
Klassifikationstoleranz ist $10^{-9}$ (FP64) bzw. $10^{-4}$ (C-Eigenwerte) relativ zu $\max|\lambda|$; die
Konsistenzprüfung lässt $2\tau$ Überlappung zu (nach dem Fluss getrennte Moden mit $|\operatorname{Re}\lambda|\le\tau$).
Der Rückfall prüft die Energiebilanz; $\operatorname{cond}(T_{11})$ wird nicht berechnet. Die Produkte
$\Psi_{\text{ref}}^{-1}\Psi$ (Substrat, RK-Scheiben, Host-Pfad der Scheiben-$S$) laufen jetzt ebenfalls über das Backend
(vorher dichte FP64-GEMM auf der CPU, bei $N=20$ die Hälfte der Laufzeit).

\subsection{Tests und Abnahme (M16)}
\begin{itemize}[nosep]
  \item Unit: Backend-Operationen in \verb|fp32| gegen \verb|fp64| (gut konditionierte Zufallsmatrizen, $n=50\dots400$):
        relativer Fehler $\le10^{-5}$; IRS-Lösen erreicht FP64-Genauigkeit ($10^{-12}$) für $\operatorname{cond}\le10^6$.
  \item Genauigkeitsstudie: alle Validierungsfälle (V1--V18) in allen drei Modi, Tabelle in \verb|docs/validation.md|
        (maximaler relativer Fehler der Effizienzen gegen \verb|fp64|, Energiebilanz, Fallback ausgelöst ja/nein).
        Zielwerte vor der Studie: \verb|mixed| $\le10^{-6}$, \verb|fp32| $\le10^{-4}$ für V1--V9. \textbf{Nach der Studie
        (v2.8):} \verb|mixed| und \verb|fp32| $\le10^{-4}$ für V1--V9 (gemessen $\le2\cdot10^{-5}$ auf CPU und RTX~3090);
        resonante 2D-Fälle (V10, V15) bis $\approx10^{-4}$. \verb|mixed| ist kaum genauer als \verb|fp32|, weil die
        Eigenprobleme (in beiden Modi C) den Fehler bestimmen; $10^{-6}$ ist mit C-Eigenwerten nicht erreichbar.
  \item Laufzeittabelle RTX~3090 und RTX 6000 Ada: 2D, $N=10,15,20$, alle Modi.
  \item \verb|AGENTS.md| Abschnitt 4: \enquote{Kern ausschließlich \texttt{complex<double>}; einfache Genauigkeit nur im
        Backend (M16).}
\end{itemize}

%=====================================================================
\section{M17: Koordinatentransformationen, adaptive räumliche Auflösung}\label{sec:ctrans}
%=====================================================================

\subsection{Formalismus}
Neue Koordinaten $X,Z$ mit
\begin{equation}
  x=f(X),\qquad z=g(Z),\qquad f(X+d_x)=f(X)+d_x,\quad g(Z+d_z)=g(Z)+d_z
  \label{eq:ct_map}
\end{equation}
(separabel, $f',g'\neq0$; ASR: reell und monoton; PML, M18: komplex in Randzonen). Die Maxwellgleichungen bleiben
unverändert, nur $\partial_x=(1/f'(X))\,\partial_X$, $\partial_z=(1/g'(Z))\,\partial_Z$. Da $f(X)-X$ periodisch ist, sind
die Felder in $X$ pseudoperiodisch mit denselben $\alpha_n$:
$e^{i\alpha_0f(X)}=e^{i\alpha_0X}\cdot(\text{periodisch})$. In der Entwicklung \eqref{eq:pseudo} nach $X,Z$ gilt
\begin{equation}
  \partial_x\to i\,A,\quad A=\tp{1/f'}\,\alpha,\qquad \partial_z\to i\,G,\quad G=\tp{1/g'}\,\gamma .
  \label{eq:ct_deriv}
\end{equation}
$\tp{1/f'}$ ist im 2D-Einzelindexraum Toeplitz in $n$ und Einheit in $m$ (umgekehrt für $\tp{1/g'}$). Deshalb
kommutieren $\tp{1/f'}$ und $\tp{1/g'}$ mit $\gamma$ bzw.\ $\alpha$ und miteinander, also $AG=GA$. $1/f'$ ist stetig;
die Laurent-Regel ist zulässig.

\begin{tcolorbox}[title=Ersetzungsregel]
In \eqref{eq:Ey}, \eqref{eq:Hy}, \eqref{eq:Mgeneral}, \eqref{eq:Mmu1a}, \eqref{eq:Mmu1b} wird jedes $\alpha$ durch $A$ und
jedes $\gamma$ durch $G$ ersetzt, \textbf{unter Beibehaltung der Reihenfolge}. Die Kurzformen in \eqref{eq:Mmu1b}
stammen aus $\alpha\,\mathcal U\,\alpha$ usw.\ mit $\mathcal U=\mathbb 1$; daraus werden
$\alpha^2\to AA$, $\alpha\gamma\to AG$, $\gamma^2\to GG$, $\gamma\alpha\to GA$. Beispiel:
$M_{14}=iA\,\mathcal E\,A-i\mathbb 1$, $M_{32}=i\big(Q_{\varepsilon,zz}-Q_{\varepsilon,zy}\mathcal EQ_{\varepsilon,yz}-AA\big)$.
Unit-Test: Für $f'\equiv g'\equiv1$ ist der Pfad mit Matrizen $A,G$ gleich dem Diagonalpfad ($10^{-15}$);
der Diagonalpfad bleibt Standard (bitgleich).
\end{tcolorbox}

\subsection{Permittivität und Normalenfeld in transformierten Koordinaten}
$\varepsilon$ wird als Funktion von $(X,Z)$ entwickelt: $\varepsilon(f(X),y,g(Z))$.
\begin{itemize}[nosep]
  \item Lamellen: Bruchstellen $X_j=f^{-1}(x_j)$, dann \eqref{eq:interval} unverändert.
  \item Polygone, Konturen, CAD-Schnitte: Rand abbilden $X(t)=f^{-1}(x(t))$, $Z(t)=g^{-1}(z(t))$ (Newton; $f$ monoton im
        physikalischen Bereich), $X'=x'/f'(X)$, dann \eqref{eq:contour} mit derselben adaptiven Quadratur.
  \item Profile: Nullstellen von $g(f(X))-y$ in $X$.
  \item Normalenfeld: \emph{physikalische} Normalen an den Bildpunkten. Die Feldkomponenten $E_x,E_z,\dots$ bleiben
        kartesische Komponenten im physikalischen Raum, daher werden die Normalen nicht transformiert.
\end{itemize}

\subsection{Homogene Halbräume, Einfall und Effizienzen}
In transformierten Koordinaten sind die Ordnungen keine Eigenvektoren mehr. Daher:
\begin{itemize}[nosep]
  \item Modenbasis von Substrat und Superstrat numerisch: Eigenzerlegung von $M$ (mit $A$, $G$, $Q_{\varepsilon}=\varepsilon\mathbb 1$),
        Klassifikation auf/ab wie bei anisotropen Medien (Abschnitt~\ref{sec:homog}).
  \item Einfallende Welle: Koeffizienten von $e^{i\alpha_0f(X)+i\gamma_0g(Z)}$ in der Basis $e^{i(\alpha_nX+\gamma_mZ)}$,
        \begin{equation}
          c_{nm}=\frac{1}{d_xd_z}\int_0^{d_x}\!\!\int_0^{d_z}e^{i\alpha_0(f(X)-X)+i\gamma_0(g(Z)-Z)}\,e^{-i(nK_xX+mK_zZ)}\,dX\,dZ
          \label{eq:ct_inc}
        \end{equation}
        (separabel, DFT), multipliziert mit der Polarisation \eqref{eq:polarization}; danach in die Modenbasis projizieren.
  \item Physikalische Rayleigh-Amplituden aus dem reflektierten bzw.\ transmittierten Tangentialfeld $U(X)=\sum_mu_me^{i\alpha_mX}$
        (1D; 2D separabel analog):
        \begin{equation}
          R_n=\frac{1}{d_x}\int_0^{d_x}U(X)\,e^{-i\alpha_nf(X)}\,f'(X)\,dX .
          \label{eq:ct_rayleigh}
        \end{equation}
        Der Integrand ist periodisch und glatt; Trapezregel mit $\ge8(2N_x+1)$ Punkten. Effizienzen dann mit
        \eqref{eq:eff_closed} aus den physikalischen Amplituden.
  \item Nahfeld: Synthese in $X$, Ausgabe auf einem Gitter in $x$ über $X=f^{-1}(x)$.
\end{itemize}

\subsection{ASR-Abbildung (Granet)}
Für Bruchstellen $0\le x_0<x_1<\dots<x_J=x_0+d_x$ (Standard: alle Materialsprünge aller Schichten in $x$, oder Liste aus
der Konfiguration) und Stärke $0\le\eta<1$, auf $[x_{j-1},x_j]$ mit $\Delta_j=x_j-x_{j-1}$:
\begin{equation}
  f(X)=X-\frac{\eta\,\Delta_j}{2\pi}\sin\frac{2\pi(X-x_{j-1})}{\Delta_j},\qquad
  f'(X)=1-\eta\cos\frac{2\pi(X-x_{j-1})}{\Delta_j}.
  \label{eq:asr}
\end{equation}
$f(x_j)=x_j$, $f'$ ist stetig (an jeder Bruchstelle $1-\eta$), monoton für $\eta<1$; ein gleichmäßiges Raster in $X$ wird an
den Bruchstellen in $x$ um den Faktor $1-\eta$ verdichtet. Lamellenbruchstellen bleiben unverändert ($X_j=x_j$).
Koeffizienten von $1/f'$ per DFT mit $\ge64(2N_x+1)$ Punkten ($1/f'$ ist stetig mit stetiger Ableitung).
Konfiguration:
{\footnotesize\begin{verbatim}
"coordinate_transform": {"x": {"type": "asr", "eta": 0.8, "breakpoints_um": "auto"}, "z": {...}}
\end{verbatim}}
Gekrümmte Ränder in 2D (\enquote{matched coordinates}, nicht separabel) sind \emph{nicht} Teil von M17.

\paragraph{Umsetzung (v2.9).} \verb|CoordinateMap| (\verb|src/geometry/coordinate_transform.cpp|) mit $f$, $f'$, $f^{-1}$
(Newton mit Bisektion). Alle Integrale über eine Periode in $X$ (Koeffizienten von $1/f'$, $C$ aus
\eqref{eq:ct_rayleigh}, $D$ aus \eqref{eq:ct_inc}, Normalenfelder der Profile) mit zusammengesetzter
Gauß-Legendre-Quadratur zwischen den Bruchstellen (16 Punkte, Panelbreite $\le d/(2\pi p_{\max})$ und
$\le$ halber Abstand der Polstelle von $1/f'$); das ersetzt DFT ($\ge64(2N+1)$ Punkte) und Trapezregel, die an den
Bruchstellen ($f'''$ springt) nur algebraisch konvergieren. Halbräume: Eigenzerlegung von $M$ (reduzierter Pfad), dann
Klassifikation aller $4n_f$ Moden mit \eqref{eq:classify} (Fluss für $|\operatorname{Re}\lambda|\le\tau$): im reduzierten
Pfad entscheidet sonst das Vorzeichen des Rundungsrauschens von $\operatorname{Im}\kappa$ propagierender Moden über
auf/ab -- in inneren Schichten unschädlich, in der Referenzbasis nicht. Homogene isotrope Medien (Halbräume und
Schichten): $PR=(A^2+G^2-\varepsilon)\otimes\mathbb 1_2$, die TE/TM-Paare sind exakt entartet und die Eigenvektoren der
vollen reduzierten Matrix (Rückwärtssubstitution der Schur-Form, nicht normal) für kleines $\eta$ schlecht konditioniert
(Energiefehler bis $3\cdot10^{-3}$ bei $\eta=10^{-6}$); daher wird nur der $n_f$-Block zerlegt und seine Eigenvektoren
für beide Komponenten verwendet. Einfall: $D$ angewandt auf das Tangentialfeld
der einfallenden Welle, Projektion auf die abwärts laufenden Moden, Anteile auf evaneszenten abwärts laufenden Moden
(Abschneiderauschen, über dem Stapel mit $e^{\lambda y}$ verstärkt) werden verworfen; Effizienzen: $C$ auf die Tangentialfelder, dann die
analytische Basis der physikalischen Ordnungen. Polygone: Rand als Kontur, abgebildet und mit \eqref{eq:contour}
integriert; Raster von Algorithmus NF bilinear auf das $X,Z$-Raster umgetastet. Profile 2D: Linien bei $z=g(Z)$, die
Bruchstellen von $g$ als zusätzliche kritische Linien. \verb|split_polarizations| wird mit Transformation nicht
genutzt. Standard-Bruchstellen \enquote{auto}: Materialsprünge der Lamellen, Ecken der Polygone.

\subsection{Tests und Abnahme (M17)}
\begin{itemize}[nosep]
  \item $\eta=0$: bitgleich mit dem Pfad ohne Transformation.
  \item V1, V2, V3, V9 mit $\eta=0.5$: Ergebnisse auf $10^{-8}$ (bei moderatem $N$) wie ohne Transformation;
        Energiebilanz und Reziprozität wie bisher. \textbf{Nach der Umsetzung (v2.9):} $10^{-8}$ gilt für die ebenen Fälle
        V1, V2 (eine Bruchstelle: $10^{-12}$; zwei: $10^{-8}$ ab $N=16$). Bei V3, V9 unterscheiden sich die Ergebnisse bei
        gleichem $N$ um den Abschneidefehler ($10^{-6}$ bei $N=40$); geprüft werden der gemeinsame Grenzwert (Abweichung
        von $N=120$ ohne Transformation kleiner als ohne ASR), die Energiebilanz ($10^{-7}$ bei $N=40$, die Umrechnung der
        Ordnungen konvergiert wie $N^{-4}$), RK $=$ RCWA und z-invariantes 2D $=$ 1D konisch ($10^{-10}$). Zusätzlich:
        transformierter Pfad mit $f'\equiv g'\equiv1$ $=$ Diagonalpfad ($10^{-12}$).
  \item \textbf{V19} (Metall-Lamellengitter TM, z.\,B.\ Gold oder Aluminium wie V13 als Lamelle): Referenz ohne ASR bei sehr
        großem $N$; mit ASR ($\eta=0.8$) gleicher Grenzwert ($10^{-4}$); Konvergenztabelle (benötigtes $N$ für $10^{-3}$
        mit/ohne ASR) in \verb|docs/validation.md|. Literaturvergleich mit Granet (1999), sofern die dortigen Daten
        bibliographisch geprüft sind.
\end{itemize}

%=====================================================================
\section{M18: PML und isolierte Strukturen}\label{sec:pml}
%=====================================================================

\subsection{Ziel}
Aperiodische (isolierte) Strukturen: Ein einzelnes Objekt wird in einer Rechenzelle mit absorbierenden Randzonen (PML)
gerechnet. Zielanwendungen (Projektinhaber): \textbf{Steg} und \textbf{Graben} (1D) sowie \textbf{Kontaktloch} (2D, Loch
in einem Film), dazu Säulen und Wellenleiterenden. Weil diese Anwendungen schrägen Einfall brauchen, gehört die
Streufeldformulierung (Abschnitt~\ref{sec:scatfield}) zum Pflichtumfang von M18. Die PML ist eine komplexe Koordinatenstreckung nach
\eqref{eq:ct_map}--\eqref{eq:ct_deriv} (M17).

\subsection{PML-Abbildung}
Rechenzelle $[0,d_x)$, physikalischer Bereich $[w,d_x-w]$ mit $f(X)=X$, Randzonen der Breite $w$ an beiden Enden.
Mit der Eindringtiefe $\xi\in[0,1]$ ($\xi=(w-X)/w$ links, $\xi=(X-d_x+w)/w$ rechts) und $u(\xi)=\sin^2(\pi\xi/2)$:
\begin{equation}
  \frac{1}{f'(X)}=\frac{1-u(\xi)}{1+i\beta\,u(\xi)} .
  \label{eq:pml}
\end{equation}
Eigenschaften: $1/f'=1$ und stetig differenzierbar am Übergang ($u'(0)=0$); $1/f'\to0$ am Zellrand (Abbildung
nach $|x|\to\infty$, vgl.\ Hugonin/Lalanne); $\operatorname{Im}f'>0$ dämpft bei der Zeitkonvention $e^{-i\omega t}$
auslaufende Wellen. Standard $w=0.15\,d_x$, $\beta=1$. \eqref{eq:pml} ist ein Vorschlag; der Agent vergleicht ihn in der
Konvergenzstudie mit der Formulierung von Hugonin/Lalanne (JOSA A 22, 1844, 2005; vor Verwendung prüfen) und
übernimmt die bessere als Standard. Analog in $z$ (2D).

\subsection{Voraussetzungen und Einfall}
\begin{itemize}[nosep]
  \item In den PML-Zonen ist jede Schicht in $x$ (bzw.\ $z$) homogen; Strukturen reichen nicht hinein. Geprüft beim
        Aufbau (Schnittkurven/Bruchstellen in den Zonen $\Rightarrow$ Fehler). Geschichtete Hintergründe (Substrat,
        Filme) laufen durch die Zonen.
  \item Einfall (a): ebene Welle bei Normaleinfall ($\alpha_0=\gamma_0=0$). Das konstante Feld ist auch in gestreckten
        Koordinaten exakt ($c_{nm}=\delta_{n0}\delta_{m0}$). Schräger Einfall ist mit seitlichen PML in der
        Totalfeldrechnung nicht darstellbar (die einfallende Welle würde in den Zonen gedämpft); dafür gibt es die
        Streufeldformulierung (c).
  \item Einfall (b): Mode eines \emph{strukturierten} Ein- oder Ausgangsmediums. Substrat und Superstrat dürfen
        $y$-invariant strukturiert sein (Lamellen, Polygon, Kontur), z.\,B.\ ein Wellenleiter, der sich in $y$ ausbreitet.
        Modenbasis aus der RCWA-Eigenzerlegung dieses Mediums; Auswahl über effektiven Index oder Modennummer.
        Normierung der Moden unkonjugiert, $\int(\mathbf E\times\mathbf H)\cdot\hat{\mathbf y}\,dx=1$ (übliche Wahl mit PML);
        für geführte Moden ist $|r|^2$ dann die Leistungsreflexion.
  \item Einfall (c): ebene Welle unter beliebigem Winkel über die Streufeldformulierung (Abschnitt~\ref{sec:scatfield}).
  \item Graben und Kontaktloch sind Löcher in Filmen, die durch die PML-Zonen laufen; das ist zulässig, weil der Film in den
        Zonen $x$-invariant ist. Hintergrund einer Schicht ist dann das Filmmaterial.
\end{itemize}

\subsection{Streufeldformulierung (schräger Einfall)}\label{sec:scatfield}
Hintergrundstapel: jede strukturierte Schicht durch ihr Hintergrundmaterial ersetzt. Seine Lösung für die einfallende
ebene Welle ist analytisch (Airy/Fresnel, bestehender Code für V2):
$\mathbf E_{\text{bg}}=e^{i(\alpha_0x+\gamma_0z)}\mathbf e(y)$, $\mathbf H_{\text{bg}}$ analog, in jeder Schicht eine Summe aus
$e^{\pm i\beta y}$-Termen. Mit $\mathbf E=\mathbf E_{\text{bg}}+\mathbf E_s$ erfüllt das Streufeld die Maxwellgleichungen
mit der Polarisationsquelle
\begin{equation}
  \mathbf D_s=\overline{\overline\varepsilon}\,\mathbf E_s+\mathbf P,\qquad
  \mathbf P=(\overline{\overline\varepsilon}-\overline{\overline\varepsilon}_{\text{bg}})\,\mathbf E_{\text{bg}},
  \label{eq:scat_source}
\end{equation}
die nur in der Struktur, also im physikalischen Bereich ($f(X)=X$), von null verschieden ist. Das Streufeld wird ohne
Blochphase entwickelt ($\alpha_n=nK_x$, $\gamma_m=mK_z$); es klingt in den PML-Zonen ab.
\begin{itemize}[nosep]
  \item \textbf{Koeffizienten der Quelle:} $\mathbf E_{\text{bg}}$ ist innerhalb einer Schicht in $x,z$ stetig, also gilt für das
        Produkt die Laurent-Regel exakt. Für jedes Material $b$ der Schicht:
        $[\chi_be^{i(\alpha_0x+\gamma_0z)}]_{pq}$ ist \eqref{eq:contour} bzw.\ \eqref{eq:polygon}/\eqref{eq:interval} mit dem
        verschobenen Wellenvektor $\mathbf k=(pK_x-\alpha_0,\ qK_z-\gamma_0)$ (der Fall $\mathbf k=0$ tritt nur für
        $\alpha_0\in K_x\mathbb Z$ auf). Damit $[\mathbf P](y)=\sum_b[\chi_be^{i(\cdot)}](\varepsilon_b-\varepsilon_{\text{bg}})\mathbf e(y)$.
  \item \textbf{Inhomogene Propagationsgleichung:} Aus \eqref{eq:Ey} mit $[D_y]=\dots+[P_y]$ und den Zeilen von \eqref{eq:ode}
        ($Q_\mu=\mathbb 1$, mit $A,G$ aus M17):
        \begin{equation}
          \frac{dF_s}{dy}=M\,F_s+\sigma(y),\qquad
          \sigma=\begin{pmatrix}-iA\,\mathcal E[P_y]\\ -iG\,\mathcal E[P_y]\\ i\big([P_z]-Q_{\varepsilon,zy}\mathcal E[P_y]\big)\\
          -i\big([P_x]-Q_{\varepsilon,xy}\mathcal E[P_y]\big)\end{pmatrix}.
          \label{eq:scat_ode}
        \end{equation}
        Vorzeichen und Faktorisierung werden mit dem Test \enquote{Streufeld + Hintergrund $=$ Totalfeld} (V20d) geprüft.
  \item \textbf{RK:} Das implizite RK-Verfahren behandelt \eqref{eq:scat_ode} direkt (Quellterm in den Stufengleichungen).
        Jede Scheibe liefert neben $S$ zwei Quellvektoren: die von der Quelle erzeugten aus- und abwärts laufenden
        Amplituden bei verschwindendem Einfall.
  \item \textbf{RCWA:} Je Scheibe partikuläre Lösung je Exponentialterm von $\sigma$: für $\sigma=\sigma_\kappa e^{i\kappa y}$ ist
        $F_p=(i\kappa\mathbb 1-M)^{-1}\sigma_\kappa e^{i\kappa y}$; liegt $i\kappa$ nahe an einem Eigenwert von $M$
        (Abstand $<10^{-8}\max|\lambda|$), wird die Scheibe mit RK gerechnet (Warnung).
  \item \textbf{Verkettung:} erweitertes Sternprodukt, das die Quellvektoren der Teilstapel mitführt (Quellen wirken wie
        zusätzliche einlaufende Amplituden an den Scheibengrenzen). Herleitung im Plan nachtragen, sobald implementiert,
        mit Test gegen die sequentielle Rechnung.
  \item \textbf{Ausgaben:} wie unten; Leistungsflüsse und Nahfeld des Streufeldes direkt, das Totalfeld als Summe.
\end{itemize}

\subsection{Ausgaben}
\begin{itemize}[nosep]
  \item Gestreutes Feld $\mathbf E_s=\mathbf E-\mathbf E_{\text{bg}}$ mit der Hintergrundlösung: derselbe Stapel, jede
        strukturierte Schicht durch ihr Hintergrundmaterial ersetzt, gleiche Transformation und gleiches $N$
        (Diskretisierungsfehler heben sich auf).
  \item Leistungsflüsse des gestreuten Feldes durch Ebenen $y=\text{const}$ über den physikalischen Bereich,
        $P=\tfrac12\operatorname{Re}\int(\mathbf E_s\times\overline{\mathbf H_s})\cdot\hat{\mathbf y}\,dx$, normiert auf die
        einfallende Intensität: Streubreite (1D, in µm) bzw.\ Streuquerschnitt (2D). Absorption aus der Flussbilanz.
  \item Bei modalem Einfall: Modenkoeffizienten $r_{jk}$, $t_{jk}$ direkt aus der S-Matrix in der Modenbasis,
        Leistungsanteile $|r_{jk}|^2$, $|t_{jk}|^2$ für geführte Moden, Abstrahlverlust $1-\sum$.
  \item Nahfeld (gesamt und gestreut) auf dem physikalischen Bereich.
\end{itemize}
Konfiguration:
{\footnotesize\begin{verbatim}
"aperiodic": {"pml": {"x": {"width_um": 0.3, "beta": 1.0}}, "excitation": "plane_wave"}   // oder {"mode": 0}
\end{verbatim}}

\paragraph{Umsetzung (v2.10).} \verb|PmlMap| (\verb|src/geometry/aperiodic.cpp|) liefert $\tp{1/f'}$ nach
\eqref{eq:pml} (Gauß-Legendre zwischen $0,w,d-w,d$); die Geometrie bleibt unverändert, weil $f(X)=X$ überall dort gilt,
wo Strukturen liegen (\verb|checkPmlZones|, auch für CAD-Schnitte und strukturierte Halbräume). Normaleinfall:
Totalfeld und Hintergrundstapel (\verb|backgroundStack|: Material am Zellrand bzw. Hintergrund, ein 1D-Profil wird
in zwei Schichten bei seiner Höhe in den Zonen geteilt) mit demselben $N$. Streufeld: statt \eqref{eq:scat_source} mit
exakten Laurent-Koeffizienten wird die Quelle als $\sigma=(M-M_{\text{bg}})F_{\text{bg}}$ gebildet ($M_{\text{bg}}$:
Hintergrundmaterial mit denselben $A,G$). Das ist zeilenweise identisch mit \eqref{eq:scat_ode}, benutzt aber für
$P=(Q_\varepsilon-\varepsilon_{\text{bg}})[\mathbf E_{\text{bg}}]$ dieselbe Faktorisierung wie das Totalfeld: Bei Li- bzw.
FFF-Regeln ist $\varepsilon\mathbf E_s$ nicht stetig, die exakten Koeffizienten von $\Delta\varepsilon\mathbf E_{\text{bg}}$
wären mit der Faktorisierung von $\mathbf E_s$ inkonsistent. Damit gilt V20d bis auf Rundung. Hintergrundfeld:
Hintergrundstapel bei $N=0$ (Airy), je Schicht vier Exponentialmoden. Schräger Einfall mit PML: das Streufeld wird
ohne Blochphase entwickelt, $\sigma$ in der Blochbasis ($\alpha_0+nK$, ohne PML-Faktor, weil $\sigma$ im physikalischen
Bereich liegt) gebildet und mit $T_{mn}=\frac1d\int_w^{d-w}e^{i(\alpha_0+(n-m)K)x}dx$ entwickelt (stetig bei
$\alpha_0\to0$). RCWA: partikuläre Lösung je Exponentialterm, aus den Werten an den Scheibengrenzen
$u=g^1_b-S_{11}g^0_b-S_{12}g^1_a$, $v=g^0_a-S_{21}g^0_b-S_{22}g^1_a$ ($g=\Psi^{-1}F_p$); Resonanz ($<10^{-8}\max|\lambda|$):
RK mit Warnung. RK: erweitertes System $(F;1)$. Verkettung: $u=u_U+U_{11}(1-L_{12}U_{21})^{-1}(u_L+L_{12}v_U)$,
$v=v_L+L_{22}(1-U_{21}L_{12})^{-1}(v_U+U_{21}u_L)$ (sequentiell und Baum), innere Amplituden analog. Ausgaben: Flüsse über das
Fenster mit dem Hintergrund als zusätzlicher \enquote{Ordnung} $(\alpha_0,\gamma_0)$; Absorption aus dem Fluss des
Totalfeldes durch die geschlossene Box (Fenster $\times$ Schichten, Seiten mit Gauß-Knoten in $y$) minus Hintergrund;
zusätzlich der seitliche Streufluss, denn nur oben $+$ unten $+$ seitlich ist von der Fenstergröße unabhängig.
Intensität: $|\mathbf S|$ der einfallenden Welle. Strukturierte Halbräume: Modenbasis aus der Eigenzerlegung von $M$;
geführte Moden: $n_{\text{eff}}=|\operatorname{Im}\lambda|$ über dem Mantelindex, $|\operatorname{Re}\lambda|<10^{-3}n_{\text{eff}}$
und $\ge90\,\%$ von $\int|E|^2$ im Fenster; Leistung über den konjugierten Fluss im Fenster (statt der unkonjugierten
Normierung, deren Integral über die gestreckte Zelle nicht konvergiert). Klassifikation der Halbraummoden: propagierende
Moden ($|\operatorname{Re}\lambda|\le10^{-6}|\lambda|$) nach dem Fluss (mit PML ist das Rauschen in
$\operatorname{Re}\lambda$ bis $10^{-7}|\lambda|$).

\subsection{Tests und Abnahme (M18)}
\begin{itemize}[nosep]
  \item \textbf{Nach der Umsetzung (v2.10):} V20a, V20d, V20e wie gefordert. V20b: Konvergenz und Unabhängigkeit von
        $w,\beta$ deutlich unter $10^{-3}$; die Absorption verlustfreier Stege ist $\le2\cdot10^{-4}$ statt $10^{-6}$ (das
        abgeschnittene PML-System erhält die Energie nur bis auf den Diskretisierungsfehler). V20c: abruptes Ende
        konvergiert auf $10^{-4}$, der Literaturvergleich ist offen (Wert nicht geprüft). V20f: Gesamtstreuung konvergiert
        auf $10^{-4}$, oben/unten nur $\sim1/L$; Reziprozität über das Winkelspektrum des Streufeldes auf einer Linie über der Struktur
        (asymmetrisch, TE/TM): $5\cdot10^{-2}$ bei $L=8\,$µm, fallend mit $L$ (Toleranz $0{,}1$; begrenzt durch die endliche Linie). V20g: Symmetrie, CAD $=$ Contour2D und
        z-invariantes 2D $=$ 1D auf $10^{-10}$; die Konvergenzstudie gehört auf das Zielsystem. Vergleich mit
        Hugonin/Lalanne nicht durchgeführt. Messwerte in \verb|docs/validation.md|.
  \item \textbf{V20a}: homogener Schichtstapel mit PML bei Normaleinfall $=$ Fresnel/Airy auf $10^{-10}$.
  \item \textbf{V20b}: isolierter dielektrischer Steg auf Substrat (1D, TE und TM): Streubreite konvergiert bei
        Zellbreite $4,6,8\times$ Stegbreite auf $10^{-3}$ und hängt nicht von $w$, $\beta$ ab ($10^{-3}$);
        Absorption $=0$ bei verlustfreiem Material ($10^{-6}$).
  \item \textbf{V20c}: Reflexion der Grundmode an einem abrupt endenden Schichtwellenleiter gegen Silberstein et~al.\
        (JOSA A 18, 2865, 2001; Wert vor Verwendung bibliographisch prüfen), Toleranz $10^{-3}$.
  \item \textbf{V20d} (Streufeldformulierung, ohne PML): periodisches Lamellengitter und \verb|Contour2D|-Säulen unter
        schrägem Einfall (auch konisch). Hintergrund $+$ Streufeld (Entwicklung mit $\alpha_n=\alpha_0+nK_x$) $=$ Totalfeldrechnung
        auf $10^{-10}$ (RCWA) bzw.\ RK-Toleranz. Damit sind \eqref{eq:scat_source}--\eqref{eq:scat_ode} unabhängig von der PML geprüft.
  \item \textbf{V20e}: Normaleinfall mit PML: Streufeldformulierung $=$ Totalfeld minus Hintergrund (Einfall a) auf $10^{-8}$.
  \item \textbf{V20f} (Graben und Steg schräg, 1D, TE/TM, $\Phi=0^\circ,30^\circ,60^\circ$): Streubreite konvergiert in der
        Zellbreite auf $10^{-3}$, unabhängig von $w,\beta$; Reziprozität der Streuamplituden bei Vertauschung von Einfalls- und
        Beobachtungsrichtung (qualitativ, Toleranz nach Studie festlegen).
  \item \textbf{V20g} (Kontaktloch, 2D, PML in $x$ und $z$): rundes Loch in einem Film auf Substrat, Normal- und Schrägeinfall;
        Konvergenz in Zellgröße und $N$ (Zielsystem); Kontaktloch als \verb|Contour2D| und als CAD-Körper identisch ($10^{-10}$).
\end{itemize}

%=====================================================================
\section{M19: Differenzierbarkeit und Optimierung}\label{sec:adjoint}
%=====================================================================

\subsection{Ziel und Umfang}
Gradienten reeller Zielfunktionen $L$ von Effizienzen oder komplexen Amplituden nach Strukturparametern, mit einem
Aufwand von höchstens 3--5 Vorwärtsrechnungen \emph{unabhängig von der Zahl der Parameter}. Parameter (Phase~1):
Dicke $y$-invarianter Schichten; Materialpermittivität (komplex); Lamellen-Bruchstellen; Polygonecken;
Kontrollpunkte und Gewichte von \verb|Contour2D|; Profilparameter (Tiefe, Verschiebung); CAD-Kontrollpunkte für
seitliche Bewegungen (Abschnitt~\ref{sec:adj_shape}). Phase~2 (optional): Einfallswinkel und Wellenlänge.

\subsection{Methode: diskrete Adjungierte}
Abgeleitet wird die \emph{tatsächlich berechnete} diskrete Abbildung (erst diskretisieren, dann differenzieren), im
Rückwärtsmodus über die aufgezeichneten Zwischenergebnisse (\verb|storeInternal| plus neue Felder). Zwei bewusste
Festlegungen, im Ergebnis als Diagnose ausgewiesen:
\begin{enumerate}[nosep]
  \item RK-Schrittfolge, Scheibenzerlegung und Bandtopologie (M14) gelten als fest.
  \item Das Normalenfeld ($\tp{N_iN_j}$) gilt als parameterunabhängig (\enquote{eingefroren}). Es beeinflusst nur die
        Konvergenzgeschwindigkeit, nicht den Grenzwert; Algorithmus~NF ist wegen der Nächster-Punkt-Suche nicht glatt.
        Der Unterschied zum vollen Differenzenquotienten wird in den Tests gemessen und fällt mit $N$.
\end{enumerate}

\subsection{Konvention (komplex, PyTorch-kompatibel)}
Für reelles $L$ und komplexe Größe $A$ heißt $\bar A:=\partial L/\partial\operatorname{Re}A+i\,\partial L/\partial\operatorname{Im}A$;
dann $dL=\operatorname{Re}\operatorname{tr}(\bar A^HdA)$. Rückwärtsregeln (Kotangenten werden \emph{addiert}):
\begin{center}
\small
\begin{tabular}{p{4.2cm}p{10.2cm}}
\toprule
Vorwärts & Rückwärts \\
\midrule
$C=AB$ & $\bar A\mathrel{+}=\bar CB^H$, $\bar B\mathrel{+}=A^H\bar C$ \\
$X=A^{-1}B$ (LU) & $\hat B=A^{-H}\bar X$ (Lösen mit $A^H$), $\bar B\mathrel{+}=\hat B$, $\bar A\mathrel{+}=-\hat BX^H$ \\
$C=A^{-1}$ & $\bar A\mathrel{+}=-C^H\bar CC^H$ \\
$z=\phi(w)$ holomorph & $\bar w\mathrel{+}=\overline{\phi'(w)}\,\bar z$ (z.\,B.\ $e^{\lambda h}$, $\sqrt{\cdot}$ für $\beta$) \\
$h$ reell, $X=\operatorname{diag}(e^{\lambda_ih})$ & $\bar h\mathrel{+}=\operatorname{Re}\sum_i\overline{\bar X_{ii}}\,\lambda_iX_{ii}$ (normiert: Faktor $k_0$ für µm) \\
$A=W\Lambda W^{-1}$ (Eig) & $\bar A=W^{-H}\big(\operatorname{diag}(\bar\lambda)+\overline{P}\circ(W^H\bar W)\big)W^H$,
  $P_{ij}=1/(\lambda_j-\lambda_i)$ ($i\neq j$), $P_{ii}=0$ \\
Toeplitz $T=\tp f$, $T_{rr'}=f_{n-n',m-m'}$ & $\bar f_{pq}\mathrel{+}=\sum_{n-n'=p,\ m-m'=q}\bar T_{rr'}$ (Summe über die Diagonalen) \\
$e=\frac{a|B_x|^2+b|B_z|^2+2c\operatorname{Re}(B_x\overline{B_z})}{D}$ \eqref{eq:eff_closed} &
  $\bar B_x\mathrel{+}=2\bar e(aB_x+cB_z)/D$, $\bar B_z\mathrel{+}=2\bar e(bB_z+cB_x)/D$ ($a,b,c,D$ reell) \\
\bottomrule
\end{tabular}
\end{center}
\textbf{Eigenzerlegung:} Die Formel setzt voraus, dass $L$ von der Normierung und Phase der Eigenvektoren nicht abhängt
(gilt für die S-Matrix). Entartete Cluster ($|\lambda_i-\lambda_j|<\tau\max|\lambda|$, $\tau=10^{-10}$) liegen in derselben
Klasse (auf/ab) und gehen nur über ihren Unterraum ein: dort $P_{ij}:=0$. Fast entartete Paare werden mit
Lorentz-Verbreiterung $P_{ij}=\overline{(\lambda_j-\lambda_i)}/(|\lambda_j-\lambda_i|^2+\delta^2)$, $\delta=10^{-12}\max|\lambda|$
behandelt. Test mit Normaleinfall auf symmetrische Gitter (exakte Entartung).

\textbf{RK-Propagator} eines Schrittes (lineare ODE, $s$ Stufen, Butcher $\mathcal A,b,c$):
$\mathbb M=\operatorname{blockdiag}(M(y+c_ih))$, $Z=\mathbb 1-h\,\mathbb M(\mathcal A\otimes\mathbb 1)$,
$K=Z^{-1}\mathbb M(\mathbf 1\otimes\mathbb 1)$, $P=\mathbb 1+h(b^T\otimes\mathbb 1)K$. Rückwärts über die Regeln oben:
$\bar K=h(b\otimes\mathbb 1)\bar P$, $\bar R=Z^{-H}\bar K$, $\bar Z=-\bar RK^H$,
$\bar{\mathbb M}\mathrel{+}=-h\bar Z(\mathcal A\otimes\mathbb 1)^H+\bar R(\mathbf 1\otimes\mathbb 1)^H$, davon die Diagonalblöcke.
(Ist der RK-Schritt im Code anders organisiert, gilt dieselbe Herleitung für die dortige Form.)

Sternprodukt \eqref{eq:star}, $T\to S$, Grenzflächen, $Q$-Matrizen (Laurent, Li-Inverse, FFF-Varianten) und der Aufbau
von $M$ \eqref{eq:Mgeneral} werden aus diesen Primitiven zusammengesetzt; jede Funktion bekommt eine eigene
Rückwärtsfunktion in \verb|src/adjoint/| und einen eigenen Gradiententest.

\subsection{Formableitungen der Koeffizienten}\label{sec:adj_shape}
Bewegt sich der Rand eines Gebiets mit der Geschwindigkeit $\mathbf V=\partial\mathbf r/\partial\theta$, dann gilt
(Reynolds/Hadamard, mit $\mathbf n\,ds=(z',-x')\,dt$), für \emph{alle} $\mathbf k$ einschließlich $\mathbf k=0$:
\begin{equation}
  \frac{\partial f_{pq}}{\partial\theta}=\frac{1}{d_xd_z}\oint\big(V_x\,z'(t)-V_z\,x'(t)\big)\,e^{-i\mathbf k\cdot\mathbf r(t)}\,dt .
  \label{eq:shape2d}
\end{equation}
NURBS-Kontur $\mathbf r(t)=\sum_jR_j(t)\mathbf P_j$: $\partial\mathbf r/\partial P_{j,x}=R_j(t)\hat{\mathbf x}$,
$\partial\mathbf r/\partial P_{j,z}=R_j(t)\hat{\mathbf z}$, für Gewichte
$\partial\mathbf r/\partial w_j=\frac{B_j(t)}{\sum_iw_iB_i(t)}(\mathbf P_j-\mathbf r(t))$. Polygonecken sind der Fall Grad~1
(Hutfunktionen). Lamellenbruchstellen ($\varepsilon=\sum_j\varepsilon_j\chi_{[x_j,x_{j+1}]}$): $\partial f_p/\partial x_j=-\frac{\varepsilon_j-\varepsilon_{j-1}}{d}\,e^{-ipKx_j}$
(Wert rechts minus links der Bruchstelle). Gleiche Quadratur wie \eqref{eq:contour}.

\textbf{CAD (3D):} Ein Flächenparameter $\theta$ (z.\,B.\ Kontrollpunkt $\mathbf P_{ij}$ in Richtung $\mathbf e$:
$\mathbf V=R_{ij}(u,v)\,\mathbf e$) verschiebt die Schnittkurve in der Ebene mit der Normalgeschwindigkeit
$v_n=(\mathbf V\cdot\mathbf N)/|\mathbf N_\parallel|$ ($\mathbf N$ Flächennormale, $\mathbf N_\parallel$ ihr Anteil in der
$(x,z)$-Ebene):
\begin{equation}
  \frac{\partial f_{pq}(y)}{\partial\theta}=\frac{1}{d_xd_z}\oint\frac{\mathbf V\cdot\mathbf N}{|\mathbf N_\parallel|}\,
  e^{-i\mathbf k\cdot\mathbf r}\,ds .
  \label{eq:shape3d}
\end{equation}
$(u,v)$ entlang der Schnittkurve aus den p-Kurven (M14). Phase~1 erlaubt nur Parameter, die keine kritische Höhe
verschieben (seitliche Bewegungen); bewegt ein Parameter eine waagrechte ebene Fläche, ist zusätzlich
$\partial y_k/\partial\theta=\mathbf V\cdot\hat{\mathbf y}$ über die Dickengradienten der beiden angrenzenden Bänder
einzurechnen; andere Fälle $\Rightarrow$ \verb|NotImplemented("M19 Phase 2")|.

\subsection{Schnittstellen}
{\footnotesize\begin{verbatim}
// C++ (include/fg/adjoint/gradient.hpp)
struct ParameterRef { std::string pointer; };      // JSON-Pointer, z. B. "/layers/2/regions/0/nurbs/points_um/3/0"
struct Objective { std::vector<std::pair<OrderRef, double>> weights; };   // L = sum w * e  (oder Kotangenten)
struct GradientResult { double value; std::vector<double> grad; Diagnostics diag; };
GradientResult gradient(const SolverOptions&, const Stack&, const Wavelength&, const Objective&,
                        const std::vector<ParameterRef>&);
// allgemeine VJP: Kotangenten fuer alle Effizienzen und/oder komplexen Amplituden -> Parameterkotangenten
std::vector<double> vjp(const Simulation&, const Cotangents&, const std::vector<ParameterRef>&);
\end{verbatim}}
Python: \verb|sim.value_and_grad(objective, params)|, \verb|sim.vjp(cotangents, params)|;
\verb|fouriergrating.torch| (\verb|torch.autograd.Function|: vorwärts Effizienzen als Tensor, rückwärts \verb|vjp|) und
\verb|fouriergrating.jax| (\verb|jax.custom_vjp| mit \verb|jax.pure_callback|); beide optional importierbar.
Komplexe Parameter (Permittivität) folgen der Konvention oben und sind damit direkt mit PyTorch kompatibel.
Speicher: Zwischenergebnisse je Scheibe; Option \verb|"checkpoint_every": k| (nur jede $k$-te Scheibe speichern, Rest neu
rechnen). GPU: Rückwärtsfunktionen über das Backend (M15/M16).

\paragraph{Umsetzung (v2.11).} Rückwärtsregeln in \verb|src/adjoint/primitives.cpp| (Tabelle oben; zusätzlich
$T\to S$, Scheiben-S-Matrix $S=HG^{-1}$ mit $\bar H=\bar SG^{-H}$, $\bar G=-S^H\bar H$, Propagationsmatrix mit Diagonal-
und dichten $A,G$, Q-Varianten Laurent/Li/FFF\_iso/FFF\_sym/FFF\_rig einschließlich der Kotangenten von $\tp{N_iN_j}$ und
$\tp{N_i}$, RK-Schritt mit Ableitung nach $h$). Der Treiber (\verb|src/solvers/gradient.cpp|) zeichnet die Kette auf
(Scheibenmatrizen, Moden, S-Matrizen, Präfixprodukte; RK: die akzeptierten Schritte, \verb|RKSteps|) und läuft
rückwärts bis zu den Koeffizientengittern. Die Schichtmodelle liefern dazu \verb|adjointSlice|: Faktorisierung,
Gitter und die Zerlegung $\varepsilon=\sum_k\varepsilon_k\chi_k$ ($1/\varepsilon$ ebenso). Formableitungen:
\eqref{eq:shape2d} mit $\mathbf V(t)=(\mathbf r_+(t)-\mathbf r_-(t))/2\delta$ aus den bei $\theta\pm\delta$ neu
erzeugten Rändern gleicher Segmentstruktur (für Kontrollpunkte exakt, $\mathbf r$ ist linear in ihnen), bei
Schnittschichten und CAD-Körpern je Scheibenhöhe aus den Schnitten (\eqref{eq:shape3d} in dieser Form);
Lamellenkanten und Profilkreuzungen analytisch. Abweichung von Festlegung 2: Das Normalenfeld der Profile ist
analytisch und glatt und wird mitdifferenziert (zentraler Differenzenquotient der Koeffizienten, ohne
$O(n^3)$-Aufwand); eingefroren wäre der Gradient bei FFF\_rig um eine Größenordnung falsch. Das Rasterfeld von
Algorithmus~NF (2D) bleibt eingefroren. Schnittstellen: \verb|fg::gradient| (JSON-Pointer), \verb|StackParameter|
(C++, beliebige Geometrie über den Stapel bei $\theta+\delta$), \verb|Simulation::valueAndGrad|/\verb|vjp|, Python
\verb|value_and_grad|/\verb|vjp|, \verb|fouriergrating.torch|, \verb|fouriergrating.jax|. Phase~2 (v2.12): Wellenlänge,
Einfallswinkel $\Phi,\Psi$, Polarisation $\delta$ und Permittivität von Superstrat (reell) und Substrat. Sie wirken über
$\alpha_n,\gamma_m$ (Kotangenten aus allen $M$ und den analytischen Basen der Halbräume, Rückwärtsregel von
\verb|isotropicModes|), über $k_0$ (normierte Dicken, $K_x=\lambda/d_x$) und über Einfallsvektor und
Effizienznormierung; diese hängen von wenigen Skalaren ab und werden mit $O(n_f)$ Aufwand per zentralem
Differenzenquotienten abgeleitet. \verb|checkpoint_every|: gespeichert werden nur jedes $k$-te Präfixprodukt, die
Scheibendaten werden rückwärts je Abschnitt neu gerechnet (Gradient bitgleich). Anisotrope Materialien: Rückwärtsregel
von \eqref{eq:QaniA}, \eqref{eq:QaniB} und Laurent bis zu den Gittern der Koeffizientenfunktionen; deren Ableitung nach
den Regionstensoren per zentralem Differenzenquotienten der Koeffizienten (holomorph, $O(n_f)$), Kanten analytisch,
Profiltiefe über das mitbewegte Normalenfeld; Lamellen, Profile und 2D-Querschnitte mit Laurent. Nicht umgesetzt:
Rückwärtsfunktionen auf dem GPU-Backend, anisotrope 2D-Querschnitte mit FFF\_B (Rasterfeld von Algorithmus~NF),
\verb|profile2d| mit anisotropen Medien, Koordinatentransformationen und aperiodische Strukturen.

\subsection{Tests und Abnahme (M19)}
\begin{itemize}[nosep]
  \item \textbf{Nach der Umsetzung (v2.11):} Rückwärtsfunktionen $10^{-6}$; Ende-zu-Ende V3, V9, V13, V15, V18
        $\le7\cdot10^{-7}$ (Laurent/Li $\le10^{-7}$), Entartung $10^{-9}$, RK $=$ RCWA $2\cdot10^{-9}$, Aufwand
        $1{,}3$--$2{,}3\times$; Messwerte in \verb|docs/validation.md|. Phase~2 (v2.12): $\lambda,\Phi,\Psi,\delta$ und
        Superstrat/Substrat gegen Differenzenquotienten $\le2\cdot10^{-7}$ (RK $2\cdot10^{-6}$, adaptive Schritte);
        anisotrope Lamellen, Profile (Laurent, FFF\_A, FFF\_B) und 2D-Säulen (Laurent) $\le3\cdot10^{-7}$.
  \item Jede Rückwärtsfunktion: zentraler Differenzenquotient mit Zufallsmatrizen ($h=10^{-7}$, Real- und Imaginärteil),
        relative Übereinstimmung $10^{-6}$.
  \item Ende-zu-Ende je Parameterart auf V3, V9, V13, V15, V18: Adjungierte gegen Differenzenquotient mit eingefrorenem
        Normalenfeld und fester Scheibenfolge $10^{-6}$; gegen den vollen Differenzenquotienten: Abweichung tabellieren.
  \item Entartung: Normaleinfall auf symmetrisches Lamellengitter, Gradient endlich und $=$ Differenzenquotient.
  \item RK- und RCWA-Gradient stimmen im Grenzwert überein (Lamellen: $10^{-8}$).
  \item Aufwand: Gradient $\le5\times$ Vorwärtsrechnung (gemessen, Bericht).
  \item Beispiel 1 (\verb|examples/python/optimize_littrow.py|): binäres Gitter, Littrow $-1$. Ordnung, Dicke und Tastgrad
        mit \verb|scipy.optimize.minimize| (L-BFGS-B); Optimum $=$ Rasterscan-Optimum auf $10^{-3}$.
  \item Beispiel 2 (\verb|examples/python/optimize_nurbs_pillar.py|): Kontrollpunkte einer \verb|Contour2D|-Säule mit
        PyTorch-Adam; Zielfunktion fällt monoton (Demonstration, kein Referenzwert).
\end{itemize}

%=====================================================================
\section{Querschnittsthemen der Version 2}\label{sec:extcross}
%=====================================================================
\begin{itemize}[nosep]
  \item \textbf{Golden Files vor dem ersten Eingriff (Beginn M12):} Für alle \verb|examples/| Effizienzen (17 signifikante
        Stellen) mit dem Stand v0.10 nach \verb|tests/regression/golden/| schreiben. Test \verb|Regression.*|: mit
        Standardoptionen auf derselben Maschine bitgleich, in CI (anderer Compiler bzw.\ andere C-Bibliothek) $10^{-11}$ bezogen auf die
        einfallende Leistung, also absolut auf den Effizienzen. Korrektur v2.3: rein relativ scheitern Ordnungen
        mit Effizienz $\approx10^{-14}$. Korrektur v2.4: $10^{-13}$ hält nur zwischen MSVC und MinGW-GCC
        ($\le1{,}2\cdot10^{-14}$); Linux-GCC~13 gegen MinGW-GCC~16 weicht bei adaptivem RK (Gold-Echelette) und
        beim 2D-Lochgitter (V10) bis $1{,}3\cdot10^{-12}$ ab (Rundung der Mathematikbibliotheken, verstärkt durch
        Schrittweitensteuerung und Kondition). Python $=$ \verb|fgsolve| gilt im selben Build bitgleich.
  \item \textbf{CI-Matrix:} \verb|cpu-release| (wie bisher), \verb|cpu-release| + \verb|FG_ENABLE_NURBFORGE=ON| (ab M14),
        Python-Wheel-Build und pytest (ab M12), ThreadSanitizer (ab M15), Doku.
  \item \textbf{Versionen:} je Meilenstein eine Minor-Version (M12 $\to$ 0.11.0, \dots, M19 $\to$ 0.18.0); Plan-Version 2.x;
        \verb|CHANGELOG.md| je Meilenstein.
  \item \textbf{Dokumentation:} \verb|docs/config_format.md| (neue Schichttypen und Blöcke), \verb|docs/user_guide.md|,
        \verb|docs/architecture.md| (Module \verb|cad|, \verb|adjoint|, Python), \verb|docs/validation.md| (V15--V20),
        \verb|docs/benchmarks.md|.
  \item \textbf{Fehlerfunde im Plan:} wie bisher (Test, Plan und Code im selben Commit, CHANGELOG).
\end{itemize}

\subsection{Entscheidungen des Projektinhabers (5.~November 2026)}
\begin{enumerate}[nosep]
  \item M18: \enquote{iso} heißt isolierte Strukturen (Steg, Graben, Kontaktloch usw.) $\Rightarrow$ Streufeldformulierung für
        schrägen Einfall ist Pflicht (Abschnitt~\ref{sec:scatfield}).
  \item Installationen: Server AlmaLinux~10 mit RTX 6000 Ada; Anwender Windows~11 mit RTX~3090 $\Rightarrow$ MSVC-Port,
        Windows-Wheels (CPU und \verb|fouriergrating-cuda12|), Server aus den Quellen (M12).
  \item CAD-Daten erzeugt ein Sprachmodell mit Nurbforge $\Rightarrow$ einzige STEP-Quelle ist der native Nurbforge-Schreiber;
        Geometrieprüfung mit Fehlercodes für die Prüfschleife des Sprachmodells (Abschnitt~\ref{sec:cadcheck}); Einheit µm.
  \item Auf der RTX~3090 ist FP32 der Standard (mit Rückfall auf FP64), auf dem Server FP64 (M16).
\end{enumerate}
Weitere Rückfragen nur, wenn ein Meilenstein auf eine im Plan nicht entschiedene Frage stößt.

\subsection*{Literatur zu Version 2 (vor Verwendung bibliographisch prüfen)}
\begin{enumerate}[label={[\arabic*]},nosep,leftmargin=2em,start=11]
  \item L.~Piegl, W.~Tiller, \emph{The NURBS Book}, 2.\ Aufl., Springer 1997.
  \item G.~Granet, \enquote{Reformulation of the lamellar grating problem through the concept of adaptive spatial resolution}, JOSA A 16, 2510 (1999).
  \item T.~Vallius, M.~Honkanen, \enquote{Reformulation of the Fourier modal method with adaptive spatial resolution: application to multilevel profiles}, Opt.\ Express 10, 24 (2002).
  \item T.~Weiss, G.~Granet, N.~A.~Gippius, S.~G.~Tikhodeev, H.~Giessen, \enquote{Matched coordinates and adaptive spatial resolution in the Fourier modal method}, Opt.\ Express 17, 8051 (2009).
  \item J.~P.~Hugonin, P.~Lalanne, \enquote{Perfectly matched layers as nonlinear coordinate transforms: a generalized formalization}, JOSA A 22, 1844 (2005).
  \item E.~Silberstein, P.~Lalanne, J.~P.~Hugonin, Q.~Cao, \enquote{Use of grating theories in integrated optics}, JOSA A 18, 2865 (2001).
  \item M.~B.~Giles, \enquote{Collected matrix derivative results for forward and reverse mode algorithmic differentiation}, in: Advances in Automatic Differentiation, Springer (2008).
  \item C.~Kim, B.~Lee, \enquote{TORCWA: GPU-accelerated Fourier modal method and gradient-based optimization for metasurface design}, Comput.\ Phys.\ Commun.\ 282, 108552 (2023).
  \item NVIDIA cuSOLVER Documentation, Iterative-Refinement-Löser (\verb|cusolverDnIRSXgesv|, \verb|cusolverDn<t1><t2>gesv|).
\end{enumerate}
